\section{德布罗意波}
\label{sec:德布罗意波}
关于波粒二象性，过去我们总是在考虑传统上被认为是波的光是否具有粒子性，但是一个相反的问题是，传统上被认为是粒子的电子，甚至是更大一些的质子和中子，是否也会具有波动性呢？在1924年德布罗意提出，\empx{任何实物粒子都兼有波和粒子的性质，都是波粒二象性的}。

我们将实物粒子对应的波称为\uwave{物质波}和\uwave{德布罗意波长}，相应的波长称为\uwave{德布罗意波长}。
\begin{OldBoxFormula}[德布罗意关系式]*
    定义微观粒子对应的物质波的频率，是能量与普朗克常数之比
    \begin{Equation}
        \nu=\frac{E}{h}=\frac{mc^2}{h}
    \end{Equation}
    定义微观粒子对应的物质波的波长，是普朗克常数与动量之比
    \begin{Equation}
        \lambda=\frac{h}{p}=\frac{h}{mv}
    \end{Equation}
\end{OldBoxFormula}
微观粒子具有波粒二象性，这需要通过实验证明，而实验确实表明，电子束也同样能呈现出干涉和衍射的波动特征。而且，电子的德布罗意波长比人类可掌控的电磁波的波长短的多，因此使用电子可以制成精度远高于光学显微镜的电子显微镜，因为波长越短，分辨本领就越强。

总而言之，现在，\empx{波和粒子之间已经不再有任何界限了}。